% ======================================================================
\section{Twierdzenie 5: problem i wynik}
% ======================================================================

\begin{frame}{Model obliczeń i problem}
  \begin{columns}[T]
    \begin{column}{0.53\textwidth}
      \begin{itemize}
        \item \textbf{Model:} word RAM ze słowami długości $O(\log N)$ bitów.
          Liczymy \emph{operacje} na liczbach całkowitych o $O(\log N)$ bitach.
        \item \textbf{Dane:} $X\in\Z^{N\times D}$, $Y\in\Z^{D\times N}$, gdzie $D\ll N$.
          Iloczyn $XY$ nazywamy \emph{cienkim}, bo macierze wejściowe są bardzo prostokątne.
        \item Zbiór $W$ pozycji $(I,J)$ macierzy $N\times N$, gdzie $I,J\in\{1,\dots,N\}$.
          Pozycje z $W$ i wpisy $XY$ na nich nazywamy \Ch{\textbf{chcianymi}}.
        \item \textbf{Cel:} policzyć tylko
          \[(XY)[I,J]=\sum_{k=1}^{D}X[I,k]\,Y[k,J]\quad\text{dla }(I,J)\in W.\]
      \end{itemize}
    \end{column}
    \begin{column}{0.45\textwidth}
      \centering
      \begin{tikzpicture}[scale=0.62]
        % X: N x D
        \draw[fill=tlozew] (0,0) rectangle (0.8,5);
        \node at (0.4,-0.35) {\small $D$};
        \node[left] at (0,2.5) {\small $N$};
        \node at (0.4,5.4) {$X$};
        % Y: D x N
        \draw[fill=tlozew] (1.4,4.2) rectangle (6.4,5);
        \node at (3.9,5.4) {$Y$};
        \node[right] at (6.4,4.6) {\small $D$};
        % XY
        \draw (1.4,0) rectangle (6.4,3.8);
        \node at (3.9,-0.35) {\small $XY$: $N\times N$};
        \foreach \x/\y in {1.9/3.2,2.7/1.1,3.3/2.6,4.6/0.6,5.8/3.4,5.2/2.0,2.2/0.5,4.0/3.0,6.0/1.3,3.6/1.6}
          \fill[chciane] (\x,\y) circle (2.2pt);
        \only<2->{
          \fill[chciane!25] (0,2.45) rectangle (0.8,2.75);
          \fill[chciane!25] (4.45,4.2) rectangle (4.75,5);
          \draw[chciane,thick] (0,2.45) rectangle (0.8,2.75);
          \draw[chciane,thick] (4.45,4.2) rectangle (4.75,5);
          \draw[chciane,dashed] (0.8,2.6) -- (4.6,2.6);
          \draw[chciane,dashed] (4.6,4.2) -- (4.6,2.6);
          \fill[chciane] (4.6,2.6) circle (3.2pt);
          \node[chciane,font=\scriptsize,align=center] at (3.2,-1.3)
            {chciany wpis = iloczyn skalarny\\ wiersza $X$ i kolumny $Y$\\ ($D$ mnożeń)};
        }
      \end{tikzpicture}
    \end{column}
  \end{columns}
\end{frame}

\begin{frame}{Główne twierdzenie rozdziału}
  \begin{twierdzenie}{5}
    Niech $D\ge4$ będzie potęgą czwórki i niech $N\ge D^{18}$. Dla danych macierzy
    $X\in\Z^{N\times D}$ i $Y\in\Z^{D\times N}$, których wpisy są liczbami całkowitymi
    o wartości bezwzględnej co najwyżej $N^{O(1)}$, oraz zbioru $W$ co najwyżej
    $N^2/\sqrt D$ pozycji macierzy $N\times N$, wpisy $(XY)[I,J]$ dla $(I,J)\in W$
    można obliczyć deterministycznie w czasie
    \[
      O\!\left(\frac{N^2\log^2 D}{D^{1/18}}\right).
    \]
  \end{twierdzenie}
  \pause
  \begin{itemize}
    \item $X$, $Y$ i $XY$ mogą być \textbf{dowolnymi gęstymi} macierzami.
    \item O zbiorze $W$ zakładamy tylko tyle, że $|W|\le N^2/\sqrt D$. Nie zakładamy żadnej struktury.
    \item Stałe (np. wykładnik $18$) dobrano tak, żeby prezentacja była prosta.
  \end{itemize}
\end{frame}

\begin{frame}{Z czym się porównujemy?}
  \centering
  \renewcommand{\arraystretch}{1.5}
  \begin{tabular}{lll}
    \toprule
    \textbf{Metoda} & \textbf{Liczba operacji} & \textbf{Uzasadnienie}\\
    \midrule
    Policzyć (wypisać) całe $XY$ & $\ge N^2$ & $XY$ ma $N^2$ wpisów\\
    \uncover<2->{Każdy chciany wpis osobno} & \uncover<2->{$|W|\cdot D\le \dfrac{N^2}{\sqrt D}\cdot D=N^2\sqrt D$} &
      \uncover<2->{$D$ mnożeń na wpis}\\
    \uncover<3->{\textbf{Twierdzenie 5}} & \uncover<3->{$O\!\left(N^2\log^2 D\cdot D^{-1/18}\right)$} & \uncover<3->{ten rozdział}\\
    \bottomrule
  \end{tabular}
  \vspace{1em}
  \uncover<4->{
  \begin{itemize}
    \item Wynik jest \textbf{wielomianowo} (w $D$) mniejszy niż \emph{wypisanie} $XY$ i niż
      $N^2/\sqrt D$ osobnych iloczynów skalarnych.
    \item Przykład: przy największym dozwolonym $D=N^{1/18}$ mamy $D^{1/18}=N^{1/324}$, czyli
      czas $O(N^{2-1/324}\log^2 N)$.
  \end{itemize}}
\end{frame}

\begin{frame}{Mapa rozdziału}
  \centering
  \begin{tikzpicture}[
      box/.style={draw,rounded corners,fill=blue!8,align=center,text width=2.7cm,
        minimum height=2.3cm,font=\small},
      arr/.style={-{Stealth[length=3mm]},thick}]
    \node[box] (s) {\textbf{2.1 Strassen}\\ rekurencja,\\ kodowanie\,/\\ dekodowanie};
    \node[box,right=0.6cm of s] (h) {\textbf{2.2 Schönhage}\\ iloczyn zewnętrzny\\ $+$ skalarny\\ $13\to10$ mnożeń};
    \node[box,right=0.6cm of h] (f) {\textbf{2.3 $\Full$}\\ $L$ poziomów:\\ wiele iloczynów\\ naraz, kafelki};
    \node[box,right=0.6cm of f,fill=red!10] (p) {\textbf{2.4 $\Pruned$}\\ \emph{nowe}: tylko\\ liście potrzebne\\ dla $W$};
    \draw[arr] (s)--(h); \draw[arr] (h)--(f); \draw[arr] (f)--(p);
  \end{tikzpicture}
  \vspace{0.6em}
  \begin{itemize}
    \item<2-> Mnożenia odbywają się w \textbf{liściach} drzewa rekursji. Dodawania na drodze w dół
      (\emph{kodowanie}) i w górę (\emph{dekodowanie}) są stosunkowo tanie.
    \item<3-> \textbf{Nowość 1:} kodowanie zależy \emph{tylko od $X$} albo \emph{tylko od $Y$},
      więc liczymy je raz i współdzielimy między wieloma wywołaniami.
    \item<4-> \textbf{Nowość 2:} przycinamy drzewo i odwiedzamy tylko liście potrzebne chcianym wpisom.
      Każdy wpis ma \emph{prywatny} liść, a pozostałe jego liście są współdzielone.
      Wychodzi $\approx N^2/D^{1/18}$ liści po $O(\log^2 D)$ operacji każdy.
  \end{itemize}
\end{frame}
